\documentclass[11pt]{article}
\usepackage{mathunal-base}
\mathunalfuente{serif}

\renewcommand{\examtitulo}{Cálculo en Varias Variables --- Segundo Parcial}
\renewcommand{\examfecha}{Semestre 01-2020, 7 de marzo de 2020}

\begin{document}

\examheader

\vspace{4pt}
\noindent Nombre: \dotfill\ Cédula: \dotfill

\vspace{4pt}
\noindent Grupo: \dotfill\ Firma: \dotfill

\vspace{4pt}
\noindent Puntaje. Solo para uso oficial:\quad 1--3 \dotfill;\quad 4 \dotfill;\quad 5 \dotfill;\quad 6 \dotfill.\quad Total: \dotfill\quad Nota: \dotfill

\vspace{6pt}
\noindent\textbf{Instrucciones.} La duración del examen es de 1 hora y 45 minutos. El examen consta de siete preguntas en dos hojas impresas por ambos lados. Antes de empezar verifique que su examen esté completo y que su nombre y datos sean asignados a Usted. No desengrape su examen ni añada otras grapas. Utilice los espacios asignados para resolver cada pregunta. No está permitido: hacer preguntas, usar calculadoras, sacar hojas en blanco, ni tomar apuntes en hojas diferentes a las del examen. No se permite que ningún celular se encuentre a la vista.

\vspace{8pt}
\noindent\textit{En cada una de las preguntas 1 a 4 marque solamente todas las afirmaciones correctas cruzando con una ``$\times$'' la caja correspondiente. Una marcación incorrecta anula una correcta. Se obtiene crédito de 6\,\% marcando todas las opciones correctas y no marcando ninguna de las incorrectas. Se obtiene crédito de 3\,\% con solo una opción mal escogida, y crédito de 0\,\% en todos los demás casos. En estas preguntas la calificación tendrá en cuenta sólo la respuesta.}

\begin{enumerate}
\item{} (6\,\%) Sea $\Omega$ la región bajo la superficie $z = 4x^2$ y encima de la región $D$ del plano $XY$ acotada por $y = 3x$, y $y = x^2$. La densidad es constante, e igual a $1$.
\begin{itemize}
\item[] \fbox{\phantom{x}}\ (a) El volumen y la masa de $\Omega$ coinciden numéricamente.
\item[] \fbox{\phantom{x}}\ (b) El volumen es $(3^5)/5$.
\item[] \fbox{\phantom{x}}\ (c) Si $(\bar x, \bar y, \bar z)$ son las coordenadas del centro de masa, entonces $3\bar x \leq \bar y \leq \bar x^2$.
\item[] \fbox{\phantom{x}}\ (d) $D$ puede escribirse como $D = \left\{(x,y) \mid 0 \leq y \leq 9,\ y/3 \leq x \leq \sqrt y\right\}$.
\end{itemize}

\item{} (6\,\%) Sea $E$ el sólido definido por $E = \left\{(x,y,z) \mid -1 \leq x \leq 1,\ x^2 \leq y \leq 1,\ 0 \leq z \leq 1 - y\right\}$. Entonces, el volumen $V$ de $E$ puede expresarse como
\begin{itemize}
\item[] \fbox{\phantom{x}}\ (a) $V = \displaystyle\int_{-1}^1 \int_0^{1-x^2} \int_{x^2}^{1-z} dy\,dz\,dx$.
\item[] \fbox{\phantom{x}}\ (b) $V = \displaystyle\int_0^1 \int_0^{1-y} \int_{-\sqrt y}^{\sqrt y} dx\,dz\,dy$.
\item[] \fbox{\phantom{x}}\ (c) $V = \displaystyle\int_0^1 \int_0^{1-z} \int_{-\sqrt y}^{\sqrt y} dx\,dy\,dz$.
\item[] \fbox{\phantom{x}}\ (d) $V = \displaystyle\int_0^1 \int_{-\sqrt y}^{\sqrt y} \int_0^{1-y} dz\,dx\,dy$.
\end{itemize}

\item{} (6\,\%) Sea $f(x,y)$ diferenciable en $\mathbb{R}^2$. Entonces:
\begin{itemize}
\item[] \fbox{\phantom{x}}\ (a) $\displaystyle\int_0^1 \int_0^x f(x,y)\,dy\,dx = \int_0^1 \int_0^y f(x,y)\,dx\,dy$.
\item[] \fbox{\phantom{x}}\ (b) $\displaystyle\int_0^1 \int_0^x f_y(x,y)\,dy\,dx = \int_0^1 f(x,x)\,dx - f(0,0)$.
\item[] \fbox{\phantom{x}}\ (c) Si $f(x,y) = -f(-x,-y)$ para $(x,y) \in D = \{(x,y) \mid x^2+y^2 \leq 9\}$, entonces $\iint_D f(x,y)\,dx\,dy = 0$.
\item[] \fbox{\phantom{x}}\ (d) Si $D = \{(x,y) \mid |x| + |y| \leq 1\}$, entonces $\iint_D f(x+y,\, 0)\,dx\,dy = \int_{-1}^1 f(u,0)\,du$. Puede considerar cambio de variables $u = x+y$, $v = x - y$.
\end{itemize}
\end{enumerate}

\vspace{4pt}
\noindent\textsc{Preguntas con procedimiento.} En las preguntas 4 a 6 responda mostrando \textit{detalladamente} el procedimiento de cálculo empleado. Indique los teoremas o procesos que sustentan sus cálculos.

\begin{enumerate}[start=4]
\item{} (22\,\%) Evalúe la integral
\[
\int_0^1 \int_{\tan^{-1} x}^{\pi/4} e^{\sec y}\,dy\,dx.
\]

\item{} (30\,\%) Sea $T(x,y) = y^2 - 2xy + x^2$ la temperatura $T$ como función de la posición $(x,y)$.
\begin{itemize}
\item[(i)] $[10\,\%]$ Usando multiplicadores de Lagrange, encuentre el sistema de ecuaciones que debe resolverse para encontrar los lugares donde se tiene la temperatura más baja y la más alta en la elipse $x^2 + 3y^2 = 1$. Numere las ecuaciones como $(1)$, $(2)$, etc.
\item[(ii)] $[20\,\%]$ Resuelva tal sistema y encuentre las temperaturas extremas absolutas en tal elipse, y todos los lugares donde se dan.
\end{itemize}

\item{} (30\,\%)
\begin{itemize}
\item[(i)] $[10\,\%]$ Sea $R$ la región en el primer cuadrante limitada por $xy = 1$, $xy = 2$, $y = x$, $y = 4x$. Calcule, usando teorema de cambio de variables,
\[
\iint_R x^2y^2\,dx\,dy.
\]
\item[(ii)] $[20\,\%]$ Encuentre la masa del sólido fuera de la esfera $x^2+y^2+z^2 = 1$, y dentro de la esfera $x^2+y^2+z^2 = 2z$, si su densidad volumétrica es $\sigma(x,y,z) = 1/(x^2+y^2+z^2)$.
\end{itemize}
\end{enumerate}

\examfooter

\end{document}
